% Marco teórico: define los conceptos técnicos y describe las ventajas de cada
% tecnología investigada (no se redacta como justificación de por qué se
% eligió ni como comparación con otras tecnologías). Estructura pedida por el
% profesor: descripción breve del marco al inicio, apartados
% numerados (1.2, 1.2.3; la numeración se activa en trabajo.tex solo para
% \solomarcoteorico), definiciones en texto corrido sin subtítulos y un
% apartado propio para la ley de protección de datos.
% Por ahora solo lo carga marco-teorico.tex (\solomarcoteorico).
% La tabla de decisiones se retiró a tablas/decisiones-marco.tex.
\section{Marco Teórico}
Este marco teórico describe los conceptos y las tecnologías en que se apoya el
sistema de acceso con códigos QR dinámicos. Primero se definen los términos
del proyecto, después se exponen las ventajas de cada tecnología y al final
se presenta la ley de protección de datos personales aplicable.

\subsection{Definiciones}
El control de acceso es el proceso mediante el cual una institución regula
quién ingresa a un espacio, sea físico o digital \parencite{alsawi2026qracs}.
El código QR es un símbolo de código de barras cuya codificación de datos,
corrección de errores y dimensiones se especifican en la norma ISO/IEC 18004
\parencite{isoiec180042015}. El HMAC es un mecanismo para autenticar
mensajes que combina una función hash criptográfica con una clave secreta
\parencite{krawczyk1997hmac}. Una contraseña de un solo uso es un valor que
solo sirve para una validación: HOTP la genera a partir de HMAC y de un
contador \parencite{mraihi2005hotp}, y TOTP es su variante basada en el
tiempo, en la que cada contraseña deja de ser válida al transcurrir un
intervalo corto \parencite{mraihi2011totp}. El ESP32 es un sistema en un chip
(SoC) con microcontrolador y conectividad Wi-Fi y Bluetooth integradas
\parencite{espressif2025esp32}, y el ESP32-CAM es un módulo de código abierto
basado en ese chip que incorpora el sensor de imagen OV2640
\parencite{nowroz2025esp32cam}.

En cuanto a las variables que se miden, el tiempo de validación es el tiempo
que transcurre desde que el estudiante llega al punto de revisión hasta que lo
cruza, sea con revisión visual o con lectura del código QR. La longitud de la
fila es el número de personas que esperan para cruzar ese punto, y su valor
máximo se registra en la franja horaria del estudio. La tasa de lecturas
fallidas es la proporción de intentos de lectura del código que el lector no
logra resolver. Son definiciones operativas del proyecto, es decir, indican
cómo se mide cada variable, y se registran con el mismo formato antes y
durante el piloto para que sean comparables. Por último, un diseño
cuasiexperimental es aquel en el que se evalúa el efecto de una intervención
sin asignar al azar a los participantes \parencite{campbellstanley1963}, y los
datos personales son cualquier información concerniente a una persona
identificada o identificable \parencite[art.~3, frac.~IX]{lgpdppso2025}.

\subsection{Control de Acceso con Código QR}
\textcite{alsawi2026qracs} señalan que los sistemas de control de acceso
protegen entornos físicos y digitales, en especial los de colegios, empresas y
áreas restringidas, y que los métodos tradicionales, como las llaves, las
tarjetas magnéticas y la biometría, presentan pérdida, duplicación, costos
altos de instalación y complejidad de mantenimiento. Cuando la validación
consiste en revisar visualmente una credencial, depende de la atención del
personal de vigilancia, y esa dependencia ya se ha documentado como fuente de
insatisfacción en otros planteles \parencite{torres2019qr}. Además, el tiempo
de espera de los estudiantes y la concentración de sus llegadas se han
estudiado como un aspecto que se puede optimizar en los servicios escolares
\parencite{talavera2025espera}. El tiempo de validación y la longitud de la
fila son las dos manifestaciones de ese problema que pueden observarse en el
sitio.

\subsubsection{Ventajas del Código QR}
Según \textcite{densowave2025qr}, las características principales del código
QR son una alta capacidad de codificación de datos, un tamaño de impresión
reducido, la resistencia a la suciedad y al daño, y la posibilidad de leerse
desde cualquier dirección; además, puede corregirse a sí mismo y recuperar los
datos originales aunque esté parcialmente manchado o dañado, y existen 40
versiones del símbolo. Su uso para el control de acceso físico se propuso
desde hace más de una década \parencite{kao2011physical}.

A estas características se suman otras ventajas para el control de acceso.
El código QR se procesa con rapidez, como se observó en transacciones móviles
\parencite{harnaningrum2024comparison}; el estudiante lo muestra desde su
propio celular, por lo que no hace falta emitir tarjetas ni registrar datos
biométricos, y la revisión de \textcite{alsawi2026qracs} lo presenta como una
alternativa flexible y de bajo costo.

\subsubsection{Código QR Dinámico}
El código QR dinámico renueva su contenido en intervalos cortos. Su ventaja
es que una captura de pantalla compartida deja de ser válida al poco tiempo,
mientras que un código estático puede copiarse y reenviarse como imagen. La
estrategia ya se probó en una universidad con un código renovado cada 20
segundos, que impidió que los estudiantes reenviaran capturas de pantalla
\parencite{nwabuwe2023dynamicqr}, y se propuso combinada con una contraseña de
un solo uso para el acceso a laboratorios universitarios
\parencite{chaabani2023labqr}. Su límite es que protege contra capturas
reutilizadas, pero no contra un estudiante que preste su celular.

\subsubsection{Firma HMAC y Contraseñas de un Solo Uso}
El código QR no incorpora cifrado propio
\parencite{harnaningrum2024comparison}, por lo que su contenido se protege con
una firma. La ventaja del HMAC es que quien no conoce la clave secreta no
puede generar una firma válida ni modificar un mensaje sin que se note
\parencite{krawczyk1997hmac}. La ventaja de TOTP es que cada valor deja de
servir al transcurrir un intervalo corto \parencite{mraihi2011totp}, lo que
impide reutilizar un código vencido. La combinación de una contraseña de un
solo uso basada en HMAC con el código QR ya se ha propuesto para autenticar
transacciones \parencite{subpratatsavee2015hmac}. Al validar, el servidor debe
tolerar un pequeño desfase entre el reloj del lector y el del celular.

\subsection{Lector del Prototipo}
El lector es el dispositivo que captura la imagen del código QR en el acceso y
la envía al servidor para su validación. Se construye con componentes de bajo
costo y conectividad inalámbrica integrada.

\subsubsection{Microcontrolador ESP32}
El ESP32 puede funcionar como un sistema autónomo o como dispositivo auxiliar
de otro microcontrolador \parencite{espressif2025esp32}. Su conectividad Wi-Fi
integrada permite que el lector se comunique con el servidor sin hardware de
red adicional.

\subsubsection{Módulo ESP32-CAM}
El ESP32-CAM es muy usado para prototipar aplicaciones de visión embebida y
destaca por su bajo costo, su versatilidad y su conectividad inalámbrica
integrada \parencite{nowroz2025esp32cam}. Estas ventajas permiten capturar la
imagen del código QR y enviarla al servidor con un solo componente. Sin
embargo, la referencia de desempeño disponible para este módulo aborda la
transmisión de video y no la decodificación de códigos QR
\parencite{nowroz2025esp32cam}, por lo que el tiempo de lectura y la tasa de
lecturas fallidas se miden con el propio prototipo, bajo distintas
condiciones de iluminación, brillo y distancia.

\subsection{Diseño Cuasiexperimental}
En la forma de preprueba y posprueba con un solo grupo, la variable de interés
se mide antes y después de la intervención en el mismo grupo
\parencite{cranmer2017pretestposttest}. Su ventaja es que permite evaluar una
intervención cuando no es posible dividir al azar a los participantes, como
ocurre con los estudiantes que llegan a un mismo acceso: basta comparar la
línea base con el piloto, como en un estudio piloto que evaluó con este diseño
un recurso con códigos QR \parencite{ortiz2026qrpilot}. Su limitación es que otros factores, como una variación en la afluencia, pueden
explicar un cambio, por lo que los resultados valen para el caso estudiado y
no se generalizan. El tipo y el alcance de la investigación se apoyan, además,
en \textcite{hernandezsampieri2014metodologia}.

\subsection{Ley de Protección de Datos Personales}
Aunque el código QR no trata datos biométricos, el sistema registra quién
ingresa y a qué hora, y esa información identifica a una persona. Por tanto,
su tratamiento queda sujeto a la legislación de protección de datos personales
en posesión de sujetos obligados, tanto general como estatal.

\subsubsection{Ley General}
La Ley General de Protección de Datos Personales en Posesión de Sujetos
Obligados establece las bases, principios y procedimientos para garantizar el
derecho de toda persona a la protección de sus datos personales en posesión de
sujetos obligados, entre los que se encuentran las autoridades, entidades,
órganos y organismos de las entidades federativas
\parencite[art.~2, fracs.~I y V]{lgpdppso2025}. Define el tratamiento como
cualquier operación o conjunto de operaciones, manuales o automatizadas,
aplicadas a los datos personales, como su obtención, registro, almacenamiento
o acceso \parencite[art.~3, frac.~XXXI]{lgpdppso2025}. De esta ley se
desprenden cuatro obligaciones para el sistema. La primera es observar los
principios de licitud, finalidad, lealtad, consentimiento, calidad,
proporcionalidad, información y responsabilidad
\parencite[art.~10]{lgpdppso2025}. La segunda es informar al titular, mediante
el aviso de privacidad, la existencia y las características principales del
tratamiento de sus datos \parencite[art.~20]{lgpdppso2025}. La tercera es
establecer y mantener medidas de seguridad de carácter administrativo, físico
y técnico que protejan los datos contra daño, pérdida, alteración, destrucción
o uso, acceso o tratamiento no autorizado \parencite[art.~25]{lgpdppso2025}.
La cuarta es permitir que el titular solicite el acceso, la rectificación, la
cancelación o la oposición al tratamiento de sus datos
\parencite[art.~37]{lgpdppso2025}.

\subsubsection{Ley del Estado de México}
En el ámbito estatal, el tratamiento se rige por la Ley de Protección de Datos
Personales en Posesión de Sujetos Obligados del Estado de México y Municipios
\parencite{edomex2017datos}. De ahí que uno de los objetivos del proyecto sea
identificar los requisitos que esta normativa impone al sistema, consultando
además a Control Escolar, antes de diseñarlo.
